% clarkedual-action-principle.tex — \input'd from proofs.tex
% Conventions: thesis/CLAUDE.md
%
% Sources: Jörn's 120-minute talk (January 2025) expanding the 1-page
%   sketch in Haim-Kislev 2017.
%
% Structure: Primal-Dual Equivalence (Clarke's Dual Action Principle for polytopes)

The proof follows \cite{HK2017} but with complete details. We first
establish a primal-dual equivalence (Clarke's Dual Action Principle
adapted to polytopes), then construct a simple minimizer via a
five-step simplification algorithm.

%%% =================================================================
%%% =================================================================


% QC: Clarke's Dual Action Principle is classical for smooth convex bodies
%   (Clarke 1979, 1981). The polytope case requires care with the gauge
%   function (not smooth) and the dual functional (integral of h_K^2).
%   Literature on this is sparse and often omits details. This writeup
%   is self-contained for a master's thesis audience.
% Downstream: establishes that minimizing action over Reeb orbits is
%   equivalent to minimizing I_K over closed curves. The dual problem
%   is easier to work with for the simplification algorithm.

We work with absolutely continuous curves
$z \in W^{1,2}([0, T], \mathbb{R}^4)$
(the Sobolev space $H^1$, which embeds into $C^0$).
% QC: W^{1,2} functions on a compact interval are absolutely continuous
%   and continuous. No need for a separate C^0 assumption.

Recall the gauge function $g_K(x) = \inf\{\lambda > 0 : x \in \lambda K\}$
(Definition~\ref{def:gauge-function}) and the support function
$h_K(v) = \max_{x \in K} \langle x, v \rangle$
(Definition~\ref{def:support-function}).
The Hamiltonian is $H = g_K^2$ (Definition~\ref{def:hamiltonian}).

\begin{lemma}[Fenchel duality]\label{lem:fenchel-duality}
% QC: Legendre-Fenchel conjugate of g_K^2 is (1/4)h_K^2.
%   The pointwise inequality is Young's inequality for this pair.
%   Equality characterizes the subgradient relation.
% Downstream: integrated along curves, this links the primal
%   (action + gauge constraint) to the dual (I_K functional).
The functions $g_K^2$ and $\tfrac{1}{4}h_K^2$ are Legendre-Fenchel conjugates:
\begin{equation}\label{eq:conjugacy}
  g_K^2(x) = \sup_y \bigl\{\langle x, y \rangle - \tfrac{1}{4} h_K^2(y)\bigr\},
  \qquad
  \tfrac{1}{4}h_K^2(y) = \sup_x \bigl\{\langle x, y \rangle - g_K^2(x)\bigr\}.
\end{equation}
Consequently, for any $x, y \in \mathbb{R}^4$:
\begin{equation}\label{eq:fenchel-inequality}
  g_K^2(x) + \tfrac{1}{4} h_K^2(y) \geq \langle x, y \rangle,
\end{equation}
with equality if and only if
$y \in \partial g_K^2(x)$ if and only if
$x \in \partial(\tfrac{1}{4} h_K^2)(y)$.
\end{lemma}

\begin{proof}
% Jörn: text approved (2f2ab3d)
\emph{Conjugacy.}
We compute the Legendre-Fenchel conjugate of~$g_K^2$ directly.
Parametrize $x = \lambda u$ with $\lambda \geq 0$ and $g_K(u) = 1$,
so $g_K^2(x) = \lambda^2$.
Then:
\[
  (g_K^2)^*(y)
  = \sup_x \bigl\{\langle x, y \rangle - g_K^2(x)\bigr\}
  = \sup_{\lambda \geq 0}\;
    \sup_{g_K(u) = 1}
    \bigl\{\lambda \langle u, y \rangle - \lambda^2\bigr\}
  = \sup_{\lambda \geq 0}
    \bigl\{\lambda\, h_K(y) - \lambda^2\bigr\},
\]
where we used $\sup_{g_K(u) = 1} \langle u, y \rangle = h_K(y)$
(the supremum of a linear functional over a convex body
is attained on its boundary).
The quadratic $\lambda h_K(y) - \lambda^2$ is maximized
at $\lambda = h_K(y)/2$, giving
\[
  (g_K^2)^*(y) = \tfrac{1}{4} h_K^2(y).
\]
This is the second relation in~\eqref{eq:conjugacy}.
The first follows by the same computation with $K$ replaced by its
polar body $K^\circ$ (whose gauge function is $h_K$ and
support function is $g_K$).

\emph{Inequality.}
By definition of the conjugate,
$(g_K^2)^*(y) \geq \langle x, y \rangle - g_K^2(x)$ for any~$x$.
Substituting $(g_K^2)^* = \tfrac{1}{4}h_K^2$ and rearranging
gives~\eqref{eq:fenchel-inequality}.

\emph{Equality condition.}
Equality holds iff $x$ attains the supremum in $(g_K^2)^*(y) = \sup_z \{\langle z, y \rangle - g_K^2(z)\}$,
i.e., $x$ maximizes $z \mapsto \langle z, y \rangle - g_K^2(z)$.
By the first-order optimality condition, this is equivalent to
$y \in \partial g_K^2(x)$.
By the same argument applied to the first relation in~\eqref{eq:conjugacy},
it is equivalent to $x \in \partial(\tfrac{1}{4}h_K^2)(y)$.
\end{proof}

\begin{definition}[Primal problem]\label{def:primal-problem}
% Jörn: text approved (05aa8b3)
% QC: restatement of the EHZ capacity minimization, using H = g_K^2.
%   The Hamiltonian formulation (differential inclusion with subdifferential)
%   is equivalent to the Reeb orbit formulation on polytopes.
The \emph{primal functional} is the action
\[
  A(\gamma) = \tfrac{1}{2}\int_0^T
  \langle -J_0\dot{\gamma},\, \gamma \rangle\, dt.
\]
The \emph{primal problem} is:
\begin{multline}\label{eq:primal-problem}
  c_{\mathrm{EHZ}}(K) = \min \Bigl\{ A(\gamma) :
  \gamma \in W^{1,2}([0,T], \mathbb{R}^4),\;
  \gamma(0) = \gamma(T),\;
  T > 0, \\
  g_K^2(\gamma) \equiv 1,\;
  -J_0\dot{\gamma} \in \partial g_K^2(\gamma) \text{ a.e.} \Bigr\}.
\end{multline}
The constraints say: $\gamma$ lies on $\partial K$
(the level set $\{g_K^2 = 1\}$) and its velocity satisfies the
Hamiltonian differential inclusion, which equals the Reeb orbit
differential inclusion (Definition~\ref{def:simple-reeb-orbit}).
The action equals the period: $A(\gamma) = T$.
\end{definition}

We now introduce an equivalent dual problem without motivation.

\begin{definition}[Dual functional and dual problem]\label{def:dual-problem}
% Jörn: text approved (05aa8b3) [revised: centering removed]
% QC: I_K depends only on the velocity dz/dt, not on z itself.
%   This is the key advantage of the dual formulation: no position constraint
%   (gamma on boundary K) and no differential inclusion constraint.
%   Translation ambiguity is harmless: I_K(z + b) = I_K(z) for any constant b.
The \emph{dual functional} is
\[
  I_K(z) = \frac{1}{4} \int_0^T h_K^2\bigl( -J_0 \dot{z}(t) \bigr) \, dt.
\]
The \emph{dual problem} is:
\begin{equation}\label{eq:dual-problem}
  \min \Bigl\{ I_K(z) :
  z \in W^{1,2}([0,T], \mathbb{R}^4),\;
  z(0) = z(T),\;
  T > 0,\;
  A(z) = T \Bigr\}.
\end{equation}
\end{definition}

\cite{HK2017} also states the dual problem without motivation
(in a slightly different form).
The dual functional and constraints can be motivated by working
the equivalence proof
(Theorem~\ref{thm:primal-dual-equivalence}) in reverse.

\begin{lemma}[Euler-Lagrange condition for dual critical points]\label{lem:dual-euler-lagrange}
% QC: Characterizes critical points via subdifferential inclusion.
%   The Lagrange multiplier is nu, the additive constant is b.
%   This is the variational condition for the constrained problem.
% Jörn: text approved (2c6faf9)
A closed curve $z \in W^{1,2}([0,T], \mathbb{R}^4)$ with $A(z) = T$ is a
critical point of the dual problem \eqref{eq:dual-problem} if and only if
there exist $\nu \in \mathbb{R}$ and $b \in \mathbb{R}^4$ such that
\begin{equation}\label{eq:dual-euler-lagrange}
  \nu z(t) + b \in \partial\bigl(\tfrac{1}{4} h_K^2\bigr)\bigl(-J_0\dot{z}(t)\bigr)
  \quad \text{a.e.,}
\end{equation}
or equivalently (by Fenchel duality),
\begin{equation}\label{eq:dual-euler-lagrange-primal}
  -J_0\dot{z}(t) \in \partial g_K^2\bigl(\nu z(t) + b\bigr)
  \quad \text{a.e.}
\end{equation}
\end{lemma}

\begin{proof}
We characterize critical points of the dual functional
$I_K(z) = \frac{1}{4}\int_0^T h_K^2(-J_0\dot{z}) \, dt$
under the constraints:
\begin{itemize}
\item closure: $z(0) = z(T)$,
\item action: $A(z) = T$.
\end{itemize}

We use the calculus of variations with Lagrange multipliers.
Form the Lagrangian incorporating the constraints:
\[
  L(z, \dot{z}) = \tfrac{1}{4}h_K^2(-J_0\dot{z}) + \alpha\langle -J_0\dot{z}, z\rangle,
\]
where $\alpha \in \mathbb{R}$ is the multiplier for the action constraint.

The Lagrangian is convex (but not differentiable) in $\dot{z}$
and smooth (linear) in~$z$.
For such Lagrangians, a critical point $z$ satisfies the
\emph{Euler--Lagrange inclusion}: there exists an absolutely continuous
$\xi(t) \in \partial_{\dot{z}} L(z(t), \dot{z}(t))$ a.e.\ such that
\begin{equation}\label{eq:euler-lagrange-inclusion-lemma}
  \dot{\xi}(t) = \nabla_z L(z(t), \dot{z}(t)) \quad\text{a.e.}
\end{equation}
(see Clarke, \emph{Optimization and Nonsmooth Analysis}, Ch.~5).

\emph{Right-hand side of~\eqref{eq:euler-lagrange-inclusion-lemma}.}
Since $\tfrac{1}{4}h_K^2(-J_0\dot{z})$ is independent of~$z$,
only the multiplier term contributes:
\[
  \nabla_z L = -\alpha J_0 \dot{z}.
\]

\emph{Left-hand side of~\eqref{eq:euler-lagrange-inclusion-lemma}.}
The subdifferential of $L$ in $\dot{z}$ is
% chain rule: ∂_ż[f(Aż)] = Aᵀ ∂f(Aż) with A = -J₀, so Aᵀ = J₀
\[
  \partial_{\dot{z}} L
  = J_0\, \partial\bigl(\tfrac{1}{4}h_K^2\bigr)(-J_0\dot{z})
    + \alpha J_0 z,
\]
using $(-J_0)^\top = J_0$ in the chain rule for $f \circ (-J_0)$.
So the Euler-Lagrange inclusion says: there exists $p(t) \in \partial(\tfrac{1}{4}h_K^2)(-J_0\dot{z}(t))$
a.e.\ such that $\xi(t) = J_0 p(t) + \alpha J_0 z(t)$ is absolutely continuous and
\[
  J_0 \dot{p} + \alpha J_0 \dot{z}
  = -\alpha J_0 \dot{z},
  \qquad\text{i.e.,}\qquad
  J_0 \dot{p} = -2\alpha J_0 \dot{z}.
\]

The equation $\dot{p} = -2\alpha \dot{z}$ a.e.\ integrates to
\[
  p(t) = c - 2\alpha\, z(t)
\]
for a constant $c \in \mathbb{R}^4$.
Since $p(t) \in \partial(\tfrac{1}{4}h_K^2)(-J_0\dot{z}(t))$:
\[
  c - 2\alpha\, z(t)
  \in \partial\bigl(\tfrac{1}{4}h_K^2\bigr)\bigl(-J_0\dot{z}(t)\bigr)
  \quad \text{a.e.}
\]
Writing $\nu = -2\alpha$ and $b = c$ gives~\eqref{eq:dual-euler-lagrange}.
\end{proof}

\begin{lemma}[Scaling and translation symmetry of dual critical points]%
\label{lem:dual-scaling-family}
% QC: The dual Euler-Lagrange equation is invariant under scaling + translation.
%   Since I_K depends only on velocity, I_K(z') = I_K(z).
%   The family parameterized by (\alpha, b) \in \mathbb{R}_{>0} \times \mathbb{R}^4
%   has constant I_K, so does not affect the minimum value of the dual problem.
% Downstream: explains why I_K(z) = T fails for general dual critical points.
%   For any Reeb orbit \gamma, we have I_K(\gamma) = T (shown in the minimum
%   equality part of Theorem 5.11).
% Jörn: text approved (9ea4716)
Let $z \in W^{1,2}([0,T], \mathbb{R}^4)$ be a closed curve that is a
critical point of the dual problem (Definition~\ref{def:dual-problem}).
For any $\alpha > 0$ and $b \in \mathbb{R}^4$, define
\begin{equation}\label{eq:scaling-family}
  z'(t) := \alpha\, z\bigl(t/\alpha^2\bigr) + b,
  \quad t \in [0, \alpha^2 T] =: [0, T'].
\end{equation}
Then $z'$ is a closed curve that is a critical point of the dual problem.
In particular, $A(z') = T'$ is satisfied. Moreover, $I_K(z') = I_K(z)$,
so this family does not affect the minimum value of the dual problem.
\end{lemma}

\begin{proof}
Closure: $z'(T') = \alpha z(T) + b = \alpha z(0) + b = z'(0)$.

Action:
\begin{align*}
  A(z')
  &= \frac{1}{2}\int_0^{T'} \langle -J_0 \dot{z}'(t), z'(t) \rangle\, dt \\
  &= \frac{1}{2}\int_0^{\alpha^2 T} \left\langle -J_0 \frac{1}{\alpha}\dot{z}(t/\alpha^2), \alpha z(t/\alpha^2) + b \right\rangle\, dt \\
  &= \frac{\alpha^2}{2}\int_0^{T} \langle -J_0 \dot{z}(s), z(s) \rangle\, ds
  \quad \text{(subst.\ $s = t/\alpha^2$, $b$ term vanishes)} \\
  &= \alpha^2 A(z) = \alpha^2 T = T'.
\end{align*}

Dual action:
\begin{align*}
  I_K(z')
  &= \frac{1}{4}\int_0^{T'} h_K^2(-J_0\dot{z}'(t))\, dt \\
  &= \frac{1}{4}\int_0^{\alpha^2 T} h_K^2\left(-J_0 \frac{1}{\alpha}\dot{z}(t/\alpha^2)\right)\, dt \\
  &= \frac{1}{4}\int_0^{\alpha^2 T} \frac{1}{\alpha^2} h_K^2(-J_0\dot{z}(t/\alpha^2))\, dt
  \quad \text{(1-homog.\ of $h_K$)} \\
  &= \frac{1}{4}\int_0^{T} h_K^2(-J_0\dot{z}(s))\, ds
  \quad \text{(subst.\ $s = t/\alpha^2$)} \\
  &= I_K(z).
\end{align*}

Critical point: By Lemma~\ref{lem:dual-euler-lagrange}, $z$ satisfies
$\nu z(s) + b_0 \in \partial(\tfrac{1}{4}h_K^2)(-J_0\dot{z}(s))$ a.e.\ for some
$\nu, b_0$. At $t$ with $s = t/\alpha^2$:
\begin{align*}
  \nu z(s) + b_0
  &\in \partial\bigl(\tfrac{1}{4}h_K^2\bigr)(-J_0\dot{z}(s)) \\
  &\Rightarrow \quad
  \frac{1}{\alpha}[\nu z(s) + b_0]
  \in \partial\bigl(\tfrac{1}{4}h_K^2\bigr)\bigl(-\tfrac{1}{\alpha}J_0\dot{z}(s)\bigr)
  \quad \text{(1-homog.\ $\partial h_K^2$)} \\
  &\Rightarrow \quad
  \frac{1}{\alpha}\left[\frac{\nu}{\alpha}(z'(t) - b) + b_0\right]
  \in \partial\bigl(\tfrac{1}{4}h_K^2\bigr)(-J_0\dot{z}'(t))
  \quad \text{(subst.\ $z(s) = (z'(t) - b)/\alpha$)} \\
  &\Rightarrow \quad
  \frac{\nu}{\alpha^2} z'(t) + \biggl(\frac{b_0}{\alpha} - \frac{\nu b}{\alpha^2}\biggr)
  \in \partial\bigl(\tfrac{1}{4}h_K^2\bigr)(-J_0\dot{z}'(t)).
\end{align*}
Thus $z'$ satisfies~\eqref{eq:dual-euler-lagrange} with $\nu' = \nu/\alpha^2$
and $b' = b_0/\alpha - \nu b/\alpha^2$.
\end{proof}
% Jörn: proof accepted as good enough (not polished)

\begin{theorem}[Primal-dual equivalence (Clarke's Dual Action Principle)]%
\label{thm:primal-dual-equivalence}
% QC: Critical points of primal and dual problems correspond via symmetry
%   classes (Lemma~\ref{lem:dual-scaling-family}). The correspondence is
%   surjective in both directions and preserves minimum values.
% Downstream: allows us to work in the dual formulation, where the
%   simplification algorithm (subsec:simplification) is natural.
% Jörn: text approved (bb06c42)
Let $K$ be a convex body. The critical points of the primal problem and of the
dual problem correspond: for a primal critical point~$\gamma$, all orbits~$z$
that are symmetry-equivalent (Lemma~\ref{lem:dual-scaling-family}) to~$\gamma$
are dual critical points. For a dual critical point~$z$, at least one
symmetry-equivalent orbit~$\gamma$ is a primal critical point.

Consequently, $\min_\gamma A(\gamma) = \min_z I_K(z) = c_{\mathrm{EHZ}}(K)$.
\end{theorem}

\begin{proof}
\emph{Forward direction: primal critical point implies symmetry class of dual critical points.}
Let $\gamma$ be a primal critical point (Hamiltonian orbit on $\partial K$)
with period~$T$. Then:
\begin{itemize}
\item $\gamma$ is closed: $\gamma(0) = \gamma(T)$,
\item $\gamma$ lies on $\partial K$: $g_K^2(\gamma(t)) \equiv 1$,
\item $\gamma$ satisfies the Hamiltonian condition:
  $-J_0\dot{\gamma} \in \partial g_K^2(\gamma)$ a.e.,
\item the action equals the period: $A(\gamma) = T$.
\end{itemize}

We show that $\gamma$ is a dual critical point.
By Fenchel duality (Lemma~\ref{lem:fenchel-duality}),
the inclusion $-J_0\dot{\gamma} \in \partial g_K^2(\gamma)$ a.e.\ is equivalent to
\[
  \gamma \in \partial\bigl(\tfrac{1}{4}h_K^2\bigr)(-J_0\dot{\gamma})
  \quad\text{a.e.}
\]
This is exactly the Euler-Lagrange condition~\eqref{eq:dual-euler-lagrange}
from Lemma~\ref{lem:dual-euler-lagrange} with $\nu = 1$ and $b = 0$.
Since $\gamma$ is closed and $A(\gamma) = T$, the curve~$\gamma$ satisfies
all constraints of the dual problem (Definition~\ref{def:dual-problem})
and is a critical point.

By Lemma~\ref{lem:dual-scaling-family}, all orbits in the symmetry class of~$\gamma$,
namely $z(t) = \alpha\gamma(t/\alpha^2) + b$ for $\alpha > 0$ and $b \in \mathbb{R}^4$,
are dual critical points.

\medskip
\emph{Reverse direction: dual critical point implies symmetry class contains primal critical point.}
Let $z$ be a dual critical point with period~$T$. By Lemma~\ref{lem:dual-euler-lagrange},
there exist $\nu \in \mathbb{R}$ and $b \in \mathbb{R}^4$ such that
\begin{equation}\label{eq:dual-el-for-z}
  \nu z(t) + b \in \partial\bigl(\tfrac{1}{4}h_K^2\bigr)(-J_0\dot{z}(t))
  \quad\text{a.e.}
\end{equation}
By Fenchel duality (Lemma~\ref{lem:fenchel-duality}), this is equivalent to
\[
  -J_0\dot{z}(t) \in \partial g_K^2(\nu z(t) + b) \quad\text{a.e.}
\]

\emph{Step 1: $\nu \neq 0$.}
Assume $\nu = 0$. Then $-J_0\dot{z}(t) \in \partial g_K^2(b)$ a.e.,
i.e., $\dot{z}(t) \in J_0\partial g_K^2(b)$ a.e.
Since $z$ is closed, $\int_0^T \dot{z} = 0$.
Since $\partial g_K^2(b)$ is convex, this implies $0 \in J_0\partial g_K^2(b)$,
hence $0 \in \partial g_K^2(b)$.
This only happens at $b = 0$.
Thus $0 \in \partial(\tfrac{1}{4}h_K^2)(\dot{z}(t))$ a.e., which implies $\dot{z} \equiv 0$,
contradicting the fact that $z$ is not constant.
Therefore $\nu \neq 0$.

\emph{Step 2: Constancy of $g_K^2(\nu z + b)$.}
The subgradient inequality gives
\[
  \frac{d}{dt} g_K^2(\nu z(t) + b)
  \geq \langle -J_0\dot{z}(t),\, \nu\dot{z}(t) \rangle
  = \nu\, \omega_0(\dot{z}(t), \dot{z}(t)) = 0
  \quad\text{a.e.}
\]
So $t \mapsto g_K^2(\nu z(t) + b)$ is non-decreasing.
Since $z$ is periodic, a non-decreasing periodic function is constant:
$g_K^2(\nu z(t) + b) \equiv c^2$ for some $c \geq 0$.
Since $\nu z + b$ is not identically zero (as $\nu \neq 0$ and $z$ is not constant),
we have $c > 0$.

\emph{Step 3: Constructing the primal critical point.}
Define
\[
  \gamma(t) := \frac{\nu z(t/\nu)}{c} + \frac{b}{c},
  \quad t \in [0, \nu T].
\]
Then $\gamma$ has period $T_\gamma = \nu T$ and
\begin{align*}
  g_K(\gamma(t))
  &= g_K\bigl(\tfrac{1}{c}(\nu z(t/\nu) + b)\bigr)
   = \tfrac{1}{c}\, g_K(\nu z(t/\nu) + b) = \tfrac{c}{c} = 1, \\
  \dot{\gamma}(t)
  &= \tfrac{\nu}{c} \cdot \tfrac{1}{\nu} \dot{z}(t/\nu)
   = \tfrac{1}{c}\dot{z}(t/\nu), \\
  -J_0\dot{\gamma}(t)
  &= -\tfrac{1}{c} J_0\dot{z}(t/\nu)
   \in \tfrac{1}{c}\, \partial g_K^2(\nu z(t/\nu) + b)
   = \partial g_K^2(\gamma(t))
  \quad\text{a.e.}
\end{align*}
We have $g_K(\gamma) \equiv 1$ and $-J_0\dot{\gamma} \in \partial g_K^2(\gamma)$ a.e.,
so $\gamma$ is a Reeb orbit (primal critical point on $\partial K$).

\emph{Step 4: Determining $\nu = c^2$ and symmetry class membership.}
For any Reeb orbit, $A(\gamma) = T_\gamma = \nu T$.
On the other hand, by direct computation with $s = t/\nu$:
\begin{align*}
  A(\gamma)
  &= \tfrac{1}{2}\!\int_0^{\nu T} \langle -J_0\dot{\gamma}(t),\, \gamma(t) \rangle\, dt
   = \tfrac{1}{2}\!\int_0^{T} \langle -\tfrac{1}{c} J_0\dot{z}(s),\,
     \tfrac{\nu z(s)}{c} + \tfrac{b}{c} \rangle\, \nu\, ds \\
  &= \tfrac{\nu}{2c^2}\!\int_0^{T} \langle -J_0\dot{z}(s),\, \nu z(s) + b \rangle\, ds
   = \tfrac{\nu^2}{c^2} A(z)
   = \tfrac{\nu^2}{c^2} T.
\end{align*}
Thus $\nu T = \nu^2 T / c^2$, hence $\nu = c^2$.

With $\nu = c^2$, we have
\[
  \gamma(t) = c\, z(t/c^2) + \frac{b}{c},
\]
which is exactly of the form $\alpha z(t/\alpha^2) + b'$ from
Lemma~\ref{lem:dual-scaling-family} with $\alpha = c > 0$ and $b' = b/c$.
Therefore $\gamma$ lies in the symmetry class of~$z$.

\medskip
\emph{Minimum equality.}
For any Reeb orbit~$\gamma$, the Fenchel equality (Lemma~\ref{lem:fenchel-duality})
at $-J_0\dot{\gamma} \in \partial g_K^2(\gamma)$ gives pointwise:
\[
  g_K^2(\gamma) + \tfrac{1}{4}h_K^2(-J_0\dot{\gamma})
  = \langle \gamma, -J_0\dot{\gamma} \rangle.
\]
Integrating over $[0,T]$ and using $g_K^2(\gamma) \equiv 1$:
\[
  T + I_K(\gamma) = 2A(\gamma) = 2T,
\]
hence $I_K(\gamma) = T = A(\gamma)$.

By Lemma~\ref{lem:dual-scaling-family}, $I_K$ is constant on symmetry classes.
We have established that every primal critical point is also a dual critical point (forward direction),
and that every dual critical point has a primal critical point in its symmetry class (reverse direction).
Therefore:
\[
  \min_z I_K(z) = \min_\gamma I_K(\gamma) = \min_\gamma A(\gamma) = c_{\mathrm{EHZ}}(K).
\]
\end{proof}

\begin{remark}\label{rem:symmetry-class-uniqueness}
% Jörn: text approved (bb06c42)
For strictly convex bodies, and for generic polytopes, exactly one Reeb orbit
lies in the symmetry class (Lemma~\ref{lem:dual-scaling-family}).
Only in degenerate cases can there be two Reeb orbits that differ by a
translation and/or spatial-temporal scaling.
\end{remark}

% [TODO: JÖRN - derive an algorithm that computes b given z]

